\chapter{Giữ Tiết Học Có Hậu Vị}
\chapterintro{Cứu lớp không chỉ là cứu giữa giờ}{Một tiết học hay là tiết giữ được dư âm sau khi đã xong.}

\section{Cuối tiết, cho lớp tự đặt tiêu đề}
\begin{tinhhuong}{Tình huống}
Muốn cả lớp nhớ tinh thần bài mà không cần lặp lại nguyên phần giảng.
\end{tinhhuong}
\begin{goiydungngay}
Hỏi: \textit{``Nếu đặt lại tên cho tiết học này sau khi đã học xong, em đặt là gì?''}
\begin{itemize}
\item Câu hỏi chốt rất gọn.
\item Học sinh buộc phải nhìn lại tinh thần bài.
\item Thường ra nhiều tên rất thú vị.
\end{itemize}
\end{goiydungngay}
\begin{cauchot}
Tiết học có hậu vị là tiết còn đọng lại thành một câu gọi tên được.
\end{cauchot}
\ngancach

\section{Một tiếng cười cuối giờ}
\begin{tinhhuong}{Tình huống}
Tiết học đã hoàn thành nội dung nhưng còn thiếu cảm giác khép lại gọn đẹp.
\end{tinhhuong}
\begin{goiydungngay}
Cho học sinh nói một câu cực ngắn kiểu: \textit{``Bài này dọa em ở chỗ nào nhất?''}
\begin{itemize}
\item Vừa vui vừa nhớ bài.
\item Không biến cuối giờ thành khảo bài nặng.
\item Rất hợp với lớp cần năng lượng nhẹ để kết thúc.
\end{itemize}
\end{goiydungngay}
\begin{cauchot}
Kết một giờ học bằng nụ cười đúng lúc thường khiến người học nhớ lâu hơn.
\end{cauchot}
\ngancach

\section{Để lại hẹn nhỏ cho tiết sau}
\begin{tinhhuong}{Tình huống}
Muốn tiết sau vào nhanh hơn, nhưng không muốn giao quá nhiều việc.
\end{tinhhuong}
\begin{goiydungngay}
Kết bằng lời hẹn rất gọn: \textit{``Tiết sau, ai mang đến cho lớp một câu trả lời thú vị nhất sẽ mở bài giúp thầy cô.''}
\begin{itemize}
\item Tạo chờ đợi nhẹ.
\item Không nặng như bài tập dài.
\item Giúp tiết sau vào nhịp nhanh hơn.
\end{itemize}
\end{goiydungngay}
\begin{cauchot}
Có hậu vị đẹp nhất là khi tiết sau đã được mở ra từ cuối tiết trước.
\end{cauchot}
\ngancach

\section{Một câu gửi lại cho lớp}
\begin{tinhhuong}{Tình huống}
Cuối tiết, thầy cô muốn để lại một dư âm nhẹ mà rõ thay vì chỉ dừng ở lời dặn.
\end{tinhhuong}
\begin{goiydungngay}
Kết bằng câu: \textit{``Nếu thầy cô gửi lại lớp một điều sau tiết này, đó là...''}
\begin{itemize}
\item Chỉ nên gửi đúng một ý.
\item Ý đó có thể là kiến thức hoặc thái độ học.
\item Nói chậm, ngắn và dứt.
\end{itemize}
\end{goiydungngay}
\begin{cauchot}
Lời gửi gọn nhưng thật thường ở lại lâu hơn lời nhắc dài.
\end{cauchot}
\ngancach

\section{Mốc nhớ cho tiết sau}
\begin{tinhhuong}{Tình huống}
Muốn bài hôm nay không tan hết trước khi tiết sau bắt đầu.
\end{tinhhuong}
\begin{goiydungngay}
Dặn lớp giữ một mốc nhớ rất nhỏ: một ví dụ, một câu, một lỗi dễ nhầm hoặc một từ khóa.
\begin{itemize}
\item Mốc nhớ càng cụ thể càng dễ giữ.
\item Tiết sau mở lại mốc đó trước khi vào bài mới.
\item Cả chuỗi bài học sẽ liền hơn.
\end{itemize}
\end{goiydungngay}
\begin{cauchot}
Hậu vị của tiết học mạnh nhất khi nó còn chỗ bám cho lần gặp sau.
\end{cauchot}